% Observed quantities are inserted by scripts/build_report.py.
\documentclass[11pt,a4paper]{article}
\usepackage[margin=1in]{geometry}
\usepackage{fontspec}
\XeTeXlinebreaklocale "zh"
\XeTeXlinebreakskip=0pt plus 1pt
\setmainfont{SimSun}[BoldFont=SimHei,ItalicFont=SimSun,BoldItalicFont=SimHei]
\newfontfamily\LatinText{Times New Roman}
\newcommand{\LatinWord}[1]{{\LatinText #1}}
\renewcommand{\abstractname}{摘要}
\renewcommand{\refname}{参考文献}
\renewcommand{\tablename}{表}
\renewcommand{\figurename}{图}
\setmonofont{Consolas}[Scale=0.9]
\usepackage{amsmath,mathtools}
\usepackage{unicode-math}
\setmathfont{Cambria Math}
\usepackage[expansion=false]{microtype}
\usepackage{graphicx,booktabs,array,tabularx,longtable,caption,enumitem,placeins}
\usepackage{flafter}
\usepackage[hidelinks]{hyperref}
\setlength{\parindent}{1em}
\setlength{\parskip}{0pt}
\setlength{\emergencystretch}{2em}
\clubpenalty=10000
\widowpenalty=10000
\setlength{\tabcolsep}{4pt}
\renewcommand{\arraystretch}{1.1}
\captionsetup{font=small,labelfont=bf}
\renewcommand{\topfraction}{0.9}
\renewcommand{\bottomfraction}{0.8}
\renewcommand{\textfraction}{0.1}
\renewcommand{\floatpagefraction}{0.75}
\setlist{itemsep=0pt,topsep=3pt}
\newcolumntype{Y}{>{\raggedright\arraybackslash}X}
\newcommand{\ResearchFigure}[3]{\begin{figure}[tbp]\centering
\includegraphics[width=\linewidth]{figures/#1.pdf}\caption{#2}\label{#3}\end{figure}}
\title{鲁棒变量选择的实验评估}
\author{}
\date{2026年10月10日}
\makeatletter
\renewcommand{\maketitle}{\begin{center}{\Large\bfseries\@title\par}\vspace{4pt}
{\small\@date\par}\end{center}\vspace{2pt}}
\makeatother
\begin{document}
\maketitle
\begin{abstract}
本报告分别验证变量组合搜索的计算开销，并整理冻结研究中的统计效用。在 16 个预先声明的确认配置上，保留完整评估已经得到的对偶分式，使完整组合求解次数减少 24.2\%，总中位运行时间变化为 +4.3\%。参考研究比较四种合成备择与两个公开数据集，变量选择、分类器拟合与测试相互独立。经验风险的全局最优并不带来一致更好的检验功效或预测。全部对照、时间上限和失败均保留，早期搜索版本的结果单独标明。
\end{abstract}

\section{实验设计与实现}\label{sec:design}
本报告区分两组证据。冻结的参考研究评估的是较早的全局割实现，其中包含五轮分数根节点求解。
它的统计结果描述当时记录的变量组合，不代表当前搜索实现的统计表现。
新的验证研究比较如何选择下一个二元组合：先在开发实例上确定规则，再在独立确认实例上检验。
两组研究使用相同的鲁棒目标，但不合并不同搜索版本或运行批次的时间结果。

实验设计沿用 Wang 等人~\cite{wang} 对计算性能、二样本检验功效和变量恢复的区分。
独立预测与小样本错误率评估参考鲁棒检验研究~\cite{xie,sinkhorn,neural}。
这些文献为评价项目提供依据；下面的实验为共同变量预算下的适配，不是对原论文表格的逐项复现。

\subsection{数据、分割与固定参数}
合成数据的类别零各坐标独立服从标准正态分布。类别一只改变一个隐藏信号块：
均值平移 $0.9$、标准差改为 $1.8$、在边缘分布不变时使用 $0.65$ 的等相关，
或向各坐标加入同一个随机正负号乘以 $1.5$。其他坐标为标准正态噪声。
坐标位置随机打乱，恢复评价仍使用稳定的原始变量身份。
相关结构变化需要联合观测；混合分布的均值相同，但形状与依赖关系不同。
\begin{table}[htbp]\centering\small\setlength{\tabcolsep}{3.5pt}
\caption{参考统计数据。选择与分类器拟合样本数均为每类数量；最终样本列注明单位。重复次数为训练种子数或重叠的真实数据分割数。}\label{tab:data}
\begin{tabularx}{\linewidth}{@{}p{25mm}rYrr@{}}
\toprule 数据 & $D/d$ & \shortstack{选择 / 拟合\\每类样本数} & \shortstack{最终\\观测} & 重复次数 \\
\midrule
合成数据 & 12/3 & 20/20 & 1000/类 & 30 \\
WDBC & 30/3 & 20/20 & 489 总数 & 20 \\
Ionosphere & 34/3 & 20/20 & 271 总数 & 20 \\
\bottomrule\end{tabularx}\end{table}
\begin{table}[htbp]\centering\small\setlength{\tabcolsep}{3.5pt}
\caption{参考优化参数。每行分别在均值与混合备择上评价。}\label{tab:computation-settings}
\begin{tabularx}{\linewidth}{@{}rrrrY@{}}
\toprule $D$ & $d$ & 源点/类 & 种子数 & 备择 \\
\midrule
8 & 2 & 8 & 10 & 均值、混合 \\
12 & 3 & 8 & 10 & 均值、混合 \\
20 & 3 & 10 & 5 & 均值、混合 \\
40 & 3 & 10 & 5 & 均值、混合 \\
80 & 3 & 10 & 5 & 均值、混合 \\
\bottomrule\end{tabularx}\end{table}
主实验两类半径均为 0.25。每个合成种子分别使用 20 个备择和 20 个零假设批次，每批每类 40 个观测。半径分析在每个机制和半径上使用 10 个种子。


在参考统计研究中，变量选择、后续模型拟合和最终观测彼此独立。中心化和样本标准差缩放只在合并的选择数据上拟合，
之后固定使用。半径、预算、正则化与带宽规则均在最终评价之前确定。
真实数据为 UCI 的 WDBC 与 Ionosphere~\cite{wdbc,ionosphere}；
特征矩阵中去除了标识符与结果标签。

\subsection{比较的变量选择方法}
六种实际可用的方法使用相同的观测预算。全局风险选择在以类别经验分布为中心的
Wasserstein 球上，最小化两类最坏半合页损失之和，运输代价使用标准化后的 $L^1$ 距离。
每个被选中的组合均完整重新拟合。有限形式、对偶割和全空间检验的推导见配套理论笔记；
本报告只讨论其实现与实验证据。
最大均值差异（MMD）通过比较类内与跨类的核相似度，衡量两组分布的差别。
高斯径向基函数（RBF）核为
$k(x,y)=\exp[-\|x-y\|_2^2/(2b^2)]$，其中 $b^2$ 是固定的平方带宽。
\begin{table}[htbp]\centering\small
\caption{预算一致的变量选择方法及固定实现参数。}\label{tab:selectors}
\begin{tabularx}{\linewidth}{@{}p{26mm}Y@{}}\toprule
方法 & 选择过程\\\midrule
全局风险选择 & 全局下界割搜索；实际可行检验提供上界。参考版本有五轮根节点求解，并以 MMD 筛选结果启动。\\
鲁棒贪心 & 逐步添加变量，每个候选都完整求解同一个鲁棒目标。若超时，保留最后完成的阶段，再用 MMD 筛选补足预算。\\
MMD 筛选 & 按单坐标无偏高斯核 MMD 排序；坐标带宽采用选择数据中跨类平方距离的中位数。\\
Wang 加性 MMD & 论文目标为加性核均值减去 $0.1$ 倍估计方差，权重满足单位范数和稀疏约束。$D\le12$ 时枚举组合并全局求解固定组合权重；更大维度采用前向选择及至多十次改进交换。\\
L1 逻辑回归 & liblinear 的 L1 拟合，$C=1$、至多 2,000 次迭代；按系数绝对值排序，并列时按选择数据方差和坐标索引处理。\\
随机选择 & 用专门且可复现的随机流，均匀选择符合预算的变量组合。\\\bottomrule
\end{tabularx}\end{table}

Wang 基线复现的是论文中的数学目标，并非逐字运行作者脚本。
所核查修订代码使用奖励方差的正号，测试过程还利用测试标签重新估计带宽。
本实现的方差项为负，所有测试带宽均由选择数据确定后固定。
其原生有符号加性统计量与共同的联合 RBF 检验分别报告。
已知信号与全部变量仅作诊断：前者在实际应用中不可获得，后者使用更大的观测预算。

每个被选组合均用相同的鲁棒分类器与 RBF SVM 评价，两者都在独立的后续数据上拟合。
鲁棒分类器使用相同半径与半合页目标；SVM 使用 $C=1$ 和只由后续拟合数据确定的中位数带宽。
这样比较的是所选坐标的效用，而不是给不同选择器更换最终模型。

\subsection{指标与不确定性}
二样本检验功效是独立备择批次中的拒绝比例；零假设拒绝率使用另外生成的、
均来自类别零分布的两个批次。共同统计量 $T$ 是联合 RBF MMD 平方的无偏估计。
带宽固定为合并选择数据平方距离的中位数，199 次固定组大小的标签置换给出
\begin{equation}\label{eq:permutation}
p=\frac{1+\#\{T_{\rm perm}\ge T_{\rm observed}-10^{-14}\}}{200},\qquad
\text{若 }p\le0.05\text{ 则拒绝}.
\end{equation}
各方法共用测试批次和置换种子。这里的零假设拒绝率与分类器在类别零上的错误率不同。

对于所选集合 $S$ 与真实信号集合 $T_*$，恢复指标为
\begin{equation}\label{eq:recovery}
\operatorname{recall}=\frac{|S\cap T_*|}{|T_*|},\qquad
\operatorname{FDP}=\frac{|S\setminus T_*|}{\max\{1,|S|\}},\qquad
\operatorname{NDP}=1-\operatorname{recall}.
\end{equation}
六种预算一致的合成选择器满足 $|S|=|T_*|=3$，因此精确率与 F1 都等于召回率，FDP 等于 NDP。
全部变量诊断不满足这个等价条件。
正文展示召回率，不重复等价指标。组合间的 Jaccard 稳定性是描述性指标，保留在完整数据中。
不为真实数据虚构已知变量真值。

令 $p_0(x)$ 为鲁棒分类器判为类别零的概率。确定性决策的阈值为 $0.5$，并列判为类别零。
其随机化错误率与共同的确定性分类指标分别为
\begin{equation}\label{eq:errors}
\begin{aligned}
\operatorname{BER}_{\rm hard}&=\tfrac12(\operatorname{FPR}+\operatorname{FNR}),\\
\operatorname{BER}_{\rm random}&=\tfrac12\{\mathbb E_0[1-p_0(X)]+\mathbb E_1[p_0(Y)]\}.
\end{aligned}
\end{equation}
FPR 是类别零被判为类别一的概率，FNR 是类别一被判为类别零的概率。
SVM 的决策阈值为零，并列时同样判为类别零。
第二个量等于半合页平均损失，并非另外一个独立结论。
AUC 使用支持类别一的评分。平衡预测交叉熵是两类负对数概率的平均，概率截断到
$10^{-12}$；它衡量概率质量，不是这里优化的半合页目标。
这些辅助指标和分类型错误率均保留在逐次实验数据中。

计算指标包括实际上下界、间隙、解出比例、完整求解过的不同组合数及累计时间。
参考研究共用一个 30 秒的选择预算，在间隙达到
$10^{-6}+10^{-5}|\mathrm{UB}|$ 时标记数值最优。
失败尝试和同模型重算均计入时间与分母。合成实验的 95\% 区间来自对训练种子的
2,000 次百分位自助重采样，测试批次先在种子内平均。
真实数据报告均值 $\pm$ 分割标准差，因为不同分割复用了数据。

\section{改进后的变量组合搜索验证}\label{sec:search}
这里的计算问题是：更充分地筛选和选择组合，能否减少昂贵的完整评估？每个选中的二元组合都完整求解，产生实际可行检验，并贡献一条切线。不将旧检验转移到新组合。

原搜索根据已完成评估预测风险均值与模型标准差，再按均值减标准差排序，候选至多八个。改进保持预测器、候选和刷新间隔不变；仅额外保留每次完整评估已经得到的安全对偶分式。筛选候选时取已有切线与分式给出的最大下界，定期主问题也保留这些分式。这里使用的是已有证书，不增加连续内层求解阶段。

实现反复执行四步：（1）排除已有下界超过当前最好风险减去一半目标容差的候选；（2）排序剩余候选；（3）完整求解选中的组合，保存其检验、切线与分式；（4）每完成三次评估，或没有候选通过筛选时，刷新全局主问题。实际检验更新上界；主问题提供全局下界并判断是否停止。

这里 $T$ 是记录的组合评估次数（求解器的 visited 计数），与满足内层精度的次数分开；受限运行的末次评估可能没有达到该精度。当前验证中的每个记录锚点都达到 $\varepsilon/4$。实际内层尝试还可能包含精度重算；连续求解次数计入其中的原始与对偶求解。主问题次数与时间则衡量整数搜索的独立成本。数值精度与全部组合枚举结果核对，而不只检查算法自身报告的外层间隙。

\subsection{开发过程与固定规则}
26 个开发实例包含两种损失，维度 6 与 8、预算 2、每类样本数 4 与 8，以及种子 3、7、11；另两个实例使用维度 12、预算 3 和每类四个样本。它们共用一个刻意选择的较差尾部坐标起点。
阶段 1 改变预测评分，或优先采用刚刷新的主问题候选。阶段 2 增加数据排序候选，或在同一点的对偶最优面上选切线；后者会额外求解一个连续问题。阶段 3 保留分式，并在每一步求解分式主问题。阶段 4 仍保留分式，但只定期刷新。单独修改评分没有减少评估次数；后两个阶段利用更强的已有下界减少评估，但每次整数主问题的成本更高。
\begin{table}[htbp]\centering\small\setlength{\tabcolsep}{3.5pt}
\caption{全部开发对照，每个阶段均为相同的 26 个配对实例。各阶段使用各自代码快照与计时批次。横线表示未记录该项测量。未完成运行保留在分母中，其已发生工作也保留。}\label{tab:search-development}
\begin{tabularx}{\linewidth}{@{}rYrrrrr@{}}
\toprule 阶段 & 搜索 & $T$ & \shortstack{连续\\求解} & 主问题 & 时间（秒） & 最优 \\
\midrule
1 & 均值减标准差 & 475 & --- & 209 & 5.188 & 26/26 \\
1 & 期望改进 & 475 & --- & 206 & 4.771 & 26/26 \\
1 & 均值减标准差，主问题优先 & 477 & --- & 218 & 4.859 & 26/26 \\
1 & 期望改进，主问题优先 & 475 & --- & 217 & 4.870 & 26/26 \\
1 & 每步切线主问题 & 467 & --- & 467 & 6.056 & 26/26 \\
2 & 原 GP 搜索 & 475 & 950 & 209 & 5.072 & 26/26 \\
2 & 增加排序候选 & 480 & 960 & 212 & 4.884 & 26/26 \\
2 & 对偶面切线 & 467 & 1419 & 206 & 17.270 & 26/26 \\
2 & 每步切线主问题 & 467 & 934 & 467 & 5.817 & 26/26 \\
3 & 原 GP 搜索 & 475 & 950 & 209 & 4.987 & 26/26 \\
3 & 每步分式主问题 & 382 & 764 & 382 & 7.722 & 26/26 \\
3 & 每步切线主问题 & 467 & 934 & 467 & 5.840 & 26/26 \\
4 & 原 GP 搜索 & 475 & 950 & 209 & 5.132 & 26/26 \\
4 & 保留分式 & 398 & 796 & 168 & 4.985 & 25/26 \\
4 & 每步分式主问题 & 382 & 764 & 382 & 7.837 & 26/26 \\
4 & 每步切线主问题 & 467 & 934 & 467 & 5.988 & 26/26 \\
\bottomrule\end{tabularx}\end{table}
阶段 4 的分式变体有一次主问题被求解器判为不可行，26 次中只有 25 次数值最优；其部分工作总数不能作为全部成功的改进结论。返回的当前检验满足已保存的下界约束。实现现在会核验这个可行见证，并在剩余预算内关闭预处理重算同一个主问题；各对照共用该数值修复。规则在确认实验前固定；下面的集成重跑只在旧开发实例上核验修复后的实现。
\begin{table}[htbp]\centering\small\setlength{\tabcolsep}{3.5pt}
\caption{数值修复后固定规则的集成重跑：26 个实例，每种方法运行一次。总数包含全部运行；这是开发核验，不是另一组确认数据。}\label{tab:search-integration}
\begin{tabularx}{\linewidth}{@{}Yrrrrrr@{}}
\toprule 搜索 & $T$ & \shortstack{连续\\求解} & 主问题 & \shortstack{主问题\\时间（秒）} & \shortstack{总时间\\（秒）} & \shortstack{最优\\运行} \\
\midrule
每步切线主问题 & 467 & 934 & 467 & 3.948 & 7.558 & 26/26 \\
原 GP 搜索 & 475 & 950 & 209 & 2.205 & 5.822 & 26/26 \\
保留分式 & 398 & 796 & 168 & 2.944 & 6.135 & 26/26 \\
\bottomrule\end{tabularx}\end{table}
\FloatBarrier
\subsection{独立确认}
16 个预先声明的确认配置使用 $(D,d)=(7,2)$ 或 $(10,3)$，每类六个样本，种子 19 与 23，半径 $(0.12,0.20)$ 或 $(0.30,0.45)$，以及半合页或逻辑损失。不同损失与半径会复用 4 份生成数据。这里直接使用生成坐标，不使用参考统计研究中的拟合标准化。类别零为高斯分布；类别一前两个坐标分别平移 1.4 与 0.9，噪声坐标的尺度从 0.3 到 3。维度与预算改变组合空间；不等半径和两种损失分别核查不对称不确定性与线性/锥求解成本。这是优化开销实验，不是新增功效实验。
各方法使用相同的数据排序起点。评分的分子是两类均值差的绝对值，分母是两类标准差与 $10^{-12}$ 三项之和。标准差使用 \texttt{ddof=0}，并列按索引稳定处理。排序不使用生成时的信号身份。
每次运行预算为 60 秒，使用一个求解线程，不评估分数根节点；绝对风险容差为 $\varepsilon=2\times10^{-5}$，相对容差为零。内层与主问题精度目标均为 $\varepsilon/4$。归一化风险预测器的先验均值为 $1/2$，协方差为 $\frac14\exp[-\|z-z'\|_1/\max(1,d)]$，正则量为 $\max\{10^{-8},(\varepsilon/(4B))^2\}$，其中半合页损失的 $B=1$，逻辑损失的 $B=2\log2$。这些参数均固定。
每个配置下各方法计时 3 次，并轮换执行顺序。下面的时间总数对每个实例先取中位数再求和；评估次数在重复间一致。16 个配置是 4 份数据上的配对设置，因此不为这些总数标注独立试验置信区间。
\begin{table}[htbp]\centering\small\setlength{\tabcolsep}{3.5pt}
\caption{确认实验的逐实例中位数总计。所有实际内层尝试均等于完整评估次数 $T$；每次评估执行两个连续求解。最优运行计数包含全部计时重复。}\label{tab:search-confirmation}
\begin{tabularx}{\linewidth}{@{}Yrrrrrr@{}}
\toprule 搜索 & $T$ & \shortstack{连续\\求解} & 主问题 & \shortstack{主问题\\时间（秒）} & \shortstack{总时间\\（秒）} & \shortstack{最优\\运行} \\
\midrule
每步切线主问题 & 601 & 1202 & 601 & 13.507 & 17.312 & 48/48 \\
原 GP 搜索 & 598 & 1196 & 314 & 8.271 & 12.715 & 48/48 \\
保留分式 & 453 & 906 & 227 & 9.610 & 13.262 & 48/48 \\
\bottomrule\end{tabularx}\end{table}
保留分式使总评估次数从 598 降到 453（减少 24.2\%），评估次数减少/相同/增加的实例分别为 15/0/1。总中位时间为 13.262 秒，原搜索为 12.715 秒（+4.3\%）；在 8/16 个实例上更快。因此，减少评估次数不能证明整体提速。分式模型的每次主问题更昂贵，尽管调用更少，主问题总时间仍增加。未改变的 GP 不能被视为这次评估减少的原因。
\ResearchFigure{search_confirmation_zh}{配对确认实例。上图：完整二元评估次数，越小越好。下图：保留分式的时间除以原 GP 时间，每种方法均使用三次运行的中位数；虚线一表示时间相同。实例 1--8 为半合页，9--16 为逻辑损失；每个损失组内依次按维度 7/10、种子 19/23、半径对排列。}{fig:search-confirmation}
全部 144/144 次确认运行均数值最优，并通过仅二元评估、证书归属与真实检验更新上界的核查。相对枚举下界参照的最大可行风险超出为 \(1.6\times10^{-7}\)。这些核查验证经验优化目标，不断言总体鲁棒性或普遍运行时间改进。


\FloatBarrier
\section{冻结参考研究：计算结果}\label{sec:compute}
各价格约束形式求解相同的二元鲁棒目标。逐项界使用 $w_{kj}\le(B/\theta_k)z_j$，
总价格界使用 $\theta_0w_0+\theta_1w_1\le Bz$，风险缩放形式用正在最小化的风险变量代替 $B$。
没有根节点求解的风险变体用于区分旧版本根节点阶段的作用。各方法均以同一个筛选组合启动。
小实例用全部组合枚举作参照，大实例考察共同时间上限下的表现。
这里 $B=1$ 是半合页损失下常数检验的两类风险和上界。
\ResearchFigure{optimization_comparison_zh}{参考版本的累计时间与解出比例。时间轴为对数刻度；
区间对种子重采样；超时与失败保留在解出比例的分母中，搜索、根节点及重算成本全部计入时间。}{fig:computation}
在 159 对已解出的约束形式/枚举结果中，最大记录可行目标差为 0.000；按合并求解容差核查，0 个区间不相容。在 $D\ge20$ 时，全局风险选择在均值变化的 15 个实例中有 15 个达到时间上限。在 $D\ge20$ 时，全局风险选择在对称混合的 15 个实例中有 15 个达到时间上限。
\begin{table}[htbp]\centering\small\setlength{\tabcolsep}{3.5pt}
\caption{$D=12$ 的参考比较。时间为累计秒数及 95\% 种子自助区间；组合数仅计不同整数子集，不含分数根节点调用。}\label{tab:reference-computation}
\begin{tabularx}{\linewidth}{@{}p{22mm}Yrrr@{}}
\toprule 备择 & 搜索 & 解出数 & 时间（秒） & 组合数 \\
\midrule
均值变化 & 逐项价格界 & 10/10 & \shortstack{4.07\\{[3.57, 4.88]}} & 220.0 \\
均值变化 & 总价格界 & 10/10 & \shortstack{4.63\\{[4.16, 5.08]}} & 219.9 \\
均值变化 & 风险（无根节点） & 10/10 & \shortstack{13.73\\{[13.05, 14.32]}} & 203.2 \\
均值变化 & 全局风险 & 10/10 & \shortstack{16.54\\{[15.74, 17.28]}} & 197.8 \\
对称混合 & 逐项价格界 & 10/10 & \shortstack{3.52\\{[3.40, 3.63]}} & 220.0 \\
对称混合 & 总价格界 & 10/10 & \shortstack{4.75\\{[4.38, 5.08]}} & 219.2 \\
对称混合 & 风险（无根节点） & 10/10 & \shortstack{13.22\\{[12.30, 14.07]}} & 200.0 \\
对称混合 & 全局风险 & 10/10 & \shortstack{16.39\\{[15.24, 17.52]}} & 197.0 \\
\bottomrule\end{tabularx}\end{table}

另行声明的完整分式比较使用不同种子，不与主时间批次合并。完整分式提供更强的值下界，但含乘积变量的整数主问题可能需要更多时间。
\begin{table}[htbp]\centering\small\setlength{\tabcolsep}{3.5pt}
\caption{较早的完整分式补充比较。末列为配对的完整分式减切线时间（秒，95\% 种子自助区间），正值表示更慢。}\label{tab:reference-ratio}
\begin{tabularx}{\linewidth}{@{}p{24mm}rrrr@{}}
\toprule 备择 & $D$ & 切线解出 & 分式解出 & \shortstack{配对时间\\差} \\
\midrule
均值变化 & 8 & 5/5 & 5/5 & \shortstack{0.19\\{[0.12, 0.25]}} \\
均值变化 & 12 & 5/5 & 3/5 & \shortstack{9.45\\{[7.38, 11.39]}} \\
均值变化 & 20 & 0/5 & 0/5 & \shortstack{0.00\\{[-0.00, 0.00]}} \\
对称混合 & 8 & 5/5 & 5/5 & \shortstack{0.31\\{[0.22, 0.40]}} \\
对称混合 & 12 & 5/5 & 4/5 & \shortstack{10.57\\{[7.36, 13.79]}} \\
对称混合 & 20 & 0/5 & 0/5 & \shortstack{0.00\\{[0.00, 0.00]}} \\
\bottomrule\end{tabularx}\end{table}


\FloatBarrier
\section{冻结参考研究：统计效用}\label{sec:statistics}
\subsection{合成检验与信号恢复}
\ResearchFigure{power_and_calibration_zh}{独立二样本功效（上排）与零假设拒绝率（下排）。
零假设轴在水平 $0.05$ 附近放大，虚线表示该水平；点与区间来自训练种子内的批次平均。只比较六种预算一致的方法。}{fig:power}
\ResearchFigure{recovery_and_prediction_zh}{信号召回率（上排，越大越好）与独立 SVM 平衡错误率
（下排，越小越好）。同列使用一个机制，同排共用指标刻度。区间对训练种子重采样。}{fig:recovery}
在均值变化上，全局选择的功效为 0.928，MMD 筛选为 0.992；配对差为 -6.3 个百分点 [-14.3, -1.3]（$n=30$）。在方差变化上，全局选择的功效为 0.360，MMD 筛选为 0.405；配对差为 -4.5 个百分点 [-19.7, 11.0]（$n=30$）。在相关结构变化上，全局选择的功效为 0.050，MMD 筛选为 0.053；配对差为 -0.3 个百分点 [-2.8, 2.0]（$n=30$）。在对称混合上，全局选择的功效为 0.770，MMD 筛选为 0.808；配对差为 -3.8 个百分点 [-20.3, 11.7]（$n=30$）。均值变化的全局搜索中有 30/30 次达到优化容差，但贪心功效为 0.965，全局选择为 0.928。经验风险最优不等于后续检验功效最大。
\begin{table}[htbp]\centering\small\setlength{\tabcolsep}{3.5pt}
\caption{Wang 所选组合的两种评价。共同功效仅使用所选坐标；原生功效还保留有符号加性权重。数值为种子均值与 95\% 区间。}\label{tab:native}
\begin{tabularx}{\linewidth}{@{}p{25mm}rrrr@{}}
\toprule 备择 & $n$ & 共同功效 & 原生功效 & 原生零假设 \\
\midrule
均值变化 & 30 & \shortstack{0.982\\{[0.962, 0.997]}} & \shortstack{0.977\\{[0.950, 0.998]}} & \shortstack{0.045\\{[0.025, 0.065]}} \\
方差变化 & 30 & \shortstack{0.328\\{[0.222, 0.435]}} & \shortstack{0.518\\{[0.408, 0.625]}} & \shortstack{0.052\\{[0.033, 0.072]}} \\
相关结构变化 & 30 & \shortstack{0.043\\{[0.025, 0.063]}} & \shortstack{0.045\\{[0.030, 0.060]}} & \shortstack{0.052\\{[0.035, 0.070]}} \\
对称混合 & 30 & \shortstack{0.763\\{[0.643, 0.872]}} & \shortstack{0.945\\{[0.915, 0.970]}} & \shortstack{0.040\\{[0.025, 0.057]}} \\
\bottomrule\end{tabularx}\end{table}
全局选择有 120/120 次独立随机化错误率高于其重新拟合的经验分类器上界的一半。这个分类器界不是变量选择目标。总体是否属于拟合的球并未得到证明，因此该超出不能称为证书失效。
\begin{table}[htbp]\centering\small\setlength{\tabcolsep}{3.5pt}
\caption{全局选择器的重新拟合分类器界与独立错误率。数值为配对种子均值；压力错误为每个种子三次预设平移最大值的平均。超出计数使用 $10^{-8}$ 容差。}\label{tab:head}
\begin{tabularx}{\linewidth}{@{}p{22mm}rrrrr@{}}
\toprule 备择 & $n$ & 拟合界 / 2 & 新样本错误 & 平移最大值 & 超出界 \\
\midrule
均值变化 & 30 & 0.251 & 0.448 & 0.462 & 30/30 \\
方差变化 & 30 & 0.275 & 0.482 & 0.484 & 30/30 \\
相关结构变化 & 30 & 0.287 & 0.497 & 0.501 & 30/30 \\
对称混合 & 30 & 0.248 & 0.423 & 0.426 & 30/30 \\
\bottomrule\end{tabularx}\end{table}


\subsection{半径与实际测量的分布平移}
半径 $0.10$、$0.25$ 和 $0.50$ 是预先设置的选项，主实验使用 $0.25$。
所选组合与重新拟合的鲁棒分类器同时改变，因此这里比较的是完整过程。
总体压力测试将各类留出观测沿后续数据均值差的共同方向移向另一类，方向的 $L^1$ 范数为一，
平移大小为 $0$、$0.25$ 和 $0.50$。最大观测错误率仅针对这三次平移，不是 Wasserstein 球上的上确界。
另一个对后续源点的平移具有显式预算可行耦合，用于直接核验经验分布证书。
\begin{table}[htbp]\centering\small\setlength{\tabcolsep}{3.5pt}
\caption{完整过程的半径敏感性。数值为种子均值与 95\% 区间；错误率为随机化平衡错误率，压力错误在每个种子内对预设平移取最大值。}\label{tab:radius}
\begin{tabularx}{\linewidth}{@{}p{22mm}rrrrr@{}}
\toprule 备择 & $\theta$ & $n$ & 功效 & 新样本错误 & 平移最大值 \\
\midrule
均值变化 & 0.10 & 10 & \shortstack{0.790\\{[0.595, 0.915]}} & \shortstack{0.494\\{[0.491, 0.497]}} & \shortstack{0.496\\{[0.493, 0.499]}} \\
均值变化 & 0.25 & 10 & \shortstack{0.855\\{[0.650, 0.980]}} & \shortstack{0.466\\{[0.450, 0.480]}} & \shortstack{0.474\\{[0.461, 0.485]}} \\
均值变化 & 0.50 & 10 & \shortstack{0.955\\{[0.915, 0.990]}} & \shortstack{0.426\\{[0.411, 0.440]}} & \shortstack{0.442\\{[0.429, 0.453]}} \\
对称混合 & 0.10 & 10 & \shortstack{0.495\\{[0.295, 0.710]}} & \shortstack{0.486\\{[0.478, 0.493]}} & \shortstack{0.487\\{[0.480, 0.493]}} \\
对称混合 & 0.25 & 10 & \shortstack{0.630\\{[0.425, 0.845]}} & \shortstack{0.446\\{[0.417, 0.469]}} & \shortstack{0.448\\{[0.419, 0.471]}} \\
对称混合 & 0.50 & 10 & \shortstack{0.970\\{[0.930, 0.995]}} & \shortstack{0.371\\{[0.352, 0.392]}} & \shortstack{0.378\\{[0.357, 0.401]}} \\
\bottomrule\end{tabularx}\end{table}
全部 1260 次预算可行源点核验中，0 次违反量超过 $10^{-7}$，最大记录非负违反量为 0.000。这个直接核验与新总体样本表现分开。


\subsection{真实数据}
\begin{table}[htbp]\centering\small\setlength{\tabcolsep}{3.5pt}
\caption{真实数据上的独立共同分类器。数值为均值 $\pm$ 重叠分割标准差。鲁棒错误为随机化 BER，SVM 错误为确定性 BER；搜索受限时使用保留的可行组合。}\label{tab:real}
\begin{tabularx}{\linewidth}{@{}p{18mm}Yrrr@{}}
\toprule 数据 & 选择器 & \shortstack{鲁棒\\随机化 BER} & SVM BER & 鲁棒 AUC \\
\midrule
WDBC & 全局风险 & 0.251\(\pm\)0.046 & 0.081\(\pm\)0.017 & 0.965\(\pm\)0.020 \\
WDBC & 鲁棒贪心 & 0.220\(\pm\)0.032 & 0.080\(\pm\)0.021 & 0.973\(\pm\)0.013 \\
WDBC & MMD 筛选 & 0.179\(\pm\)0.013 & 0.084\(\pm\)0.017 & 0.973\(\pm\)0.011 \\
WDBC & Wang 加性 MMD & 0.180\(\pm\)0.013 & 0.082\(\pm\)0.014 & 0.975\(\pm\)0.010 \\
WDBC & L1 逻辑回归 & 0.230\(\pm\)0.034 & 0.078\(\pm\)0.020 & 0.971\(\pm\)0.017 \\
WDBC & 随机 & 0.310\(\pm\)0.066 & 0.156\(\pm\)0.078 & 0.893\(\pm\)0.093 \\
Ionosphere & 全局风险 & 0.300\(\pm\)0.045 & 0.251\(\pm\)0.094 & 0.869\(\pm\)0.053 \\
Ionosphere & 鲁棒贪心 & 0.285\(\pm\)0.035 & 0.202\(\pm\)0.081 & 0.885\(\pm\)0.045 \\
Ionosphere & MMD 筛选 & 0.284\(\pm\)0.044 & 0.195\(\pm\)0.053 & 0.865\(\pm\)0.055 \\
Ionosphere & Wang 加性 MMD & 0.273\(\pm\)0.042 & 0.181\(\pm\)0.058 & 0.877\(\pm\)0.051 \\
Ionosphere & L1 逻辑回归 & 0.284\(\pm\)0.034 & 0.198\(\pm\)0.049 & 0.880\(\pm\)0.046 \\
Ionosphere & 随机 & 0.330\(\pm\)0.028 & 0.303\(\pm\)0.069 & 0.835\(\pm\)0.034 \\
\bottomrule\end{tabularx}\end{table}
全局搜索在 40/40 次分割中达到时间上限，因此不能将这些组合称为已认证最优解。


\FloatBarrier
\section{运行记录与可复现性}\label{sec:records}
\begin{table}[htbp]\centering\small\setlength{\tabcolsep}{3.5pt}
\caption{保留的全部冻结运行结果。最优指在声明容差下解出鲁棒目标。其他包含基线与诊断完成，不代表鲁棒全局最优。}\label{tab:outcomes}
\begin{tabularx}{\linewidth}{@{}Yrrrrr@{}}
\toprule 组别 & 记录 & 最优 & 时间上限 & 失败 & 其他 \\
\midrule
合成统计 & 960 & 114 & 6 & 0 & 840 \\
优化 & 320 & 199 & 120 & 1 & 0 \\
真实数据 & 240 & 0 & 40 & 0 & 200 \\
半径 & 60 & 58 & 0 & 2 & 0 \\
分式补充 & 60 & 37 & 23 & 0 & 0 \\
\bottomrule\end{tabularx}\end{table}
最初 12 次失败在原选择预算的剩余部分内重算：最终 5 次最优、6 次达到时间上限，1 次失败。最终记录共有 3 次失败，包含后来半径实验的结果。原尝试与重算成本全部保留。
参考计算使用 AMD Ryzen 9 7945HX with Radeon Graphics、一个求解线程，以及 SciPy 内置的 HiGHS 1.12.0（SciPy 1.18.1）。逻辑损失核验使用 CVXPY 1.9.3 和 Clarabel 0.11.1。累计选择时间与全部流程时间分别为 2.94 和 2.96 小时，包含未成功尝试；并发运行时间的总和不是实际日历耗时。


完整冻结记录、所有辅助指标、单次所选组合、预处理和优化轨迹保留在
\texttt{results/research\_20261003/per\_trial.csv} 及相应 JSON 记录中。
新开发与确认结果单独冻结在 \texttt{reports/verification/20261010/next\_selection} 下。
报告清单保存输入与图的哈希；重新生成时会把原始记录与配置哈希、逐次数据表进行核对。
优化精度和拟合得到的经验球证书都不能证明总体覆盖。
准则与独立统计效用之间实际观察到的差距也作为结果保留。

\FloatBarrier
\begin{thebibliography}{9}
\bibitem{wang} J. Wang, S. S. Dey, and Y. Xie.
\emph{Variable Selection for Kernel Two-Sample Tests}. 2024.
\href{https://arxiv.org/abs/2302.07415}{arXiv:2302.07415v4}.
代码核查版本为
\href{https://github.com/WalterBabyRudin/MMDVar_Selection}{MMDVar\_Selection}
的提交 \texttt{031f740e9ca6}。
\bibitem{xie} L. Xie, R. Gao, and Y. Xie.
\emph{Robust Hypothesis Testing with Wasserstein Uncertainty Sets}. 2021.
\href{https://arxiv.org/abs/2105.14348}{arXiv:2105.14348}.
\bibitem{sinkhorn} J. Wang and Y. Xie.
\emph{A Data-Driven Approach to Robust Hypothesis Testing Using Sinkhorn Uncertainty Sets}.
\bibitem{neural} F. Zhang, T. Liu, and J. Wang.
\emph{Generative Neural Networks for Sinkhorn Distributionally Robust Hypothesis Testing}. 2026.
\href{https://github.com/Arthas819/SDRHT}{Author code and experiment repository}.
\bibitem{wdbc} UCI Machine Learning Repository.
\emph{Breast Cancer Wisconsin (Diagnostic)}. \href{https://doi.org/10.24432/C5DW2B}{DOI:10.24432/C5DW2B}.
\bibitem{ionosphere} UCI Machine Learning Repository.
\emph{Ionosphere}. \href{https://doi.org/10.24432/C5W01B}{DOI:10.24432/C5W01B}.
\end{thebibliography}
\end{document}
